% INSERT AFTER: online_appendix.tex lines 706-707 ("\subsection{A.6. Smooth noise, atomless costs, and primitive" to "nonemptiness}\label{oa-a-extensions}").
% INSERT BEFORE: "For logistic noise," at line 709. Keep the existing subsection and equation numbers.
% NEW LABELS: oa-a-laplace-extreme, result-lemma-oa-laplace-extreme.
% NEW SYMBOLS: mathcal C_L, L, ell_b (Laplace density), S (survival function),
% beta_f (ROC curve), varphi (statistical decision rule),
% h_varphi, Phi, v (ODE objects), Y_n (independent Laplace variables),
% W (extra noise), and chi_log (the logistic characteristic function).
% Existing pi is the class prior throughout; the characteristic function uses its product form.
% NUMBERS: none. No new bibliography record.
% SOURCES: Q2_laplace_technical R8-R10; the general-prior weights use Q1_N_main Claim 6.

\subsubsection{Laplace is the extreme member of the class}\label{oa-a-laplace-extreme}

Write \(\mathcal C_L\) for the positive, absolutely continuous densities
with \(|f'|\le Lf\) almost everywhere.
This class has no location restriction.
Part (i) imposes median zero only for the tail comparison.
\par
Set \(b=1/L\) and \(\ell_b(z)=e^{-|z|/b}/(2b)\).
The score bound gives \(\int|f'|\le L\).
The convolution argument in A.3 therefore applies.
The three claims below hold for every \(L>0\).
The equilibrium result requires \(L<1\) in addition.
\par
All experiments use the same prior \(0<\pi<1\) and the same order profile.
The price comparison also fixes the auction primitives.
Let \(\beta_f(\alpha)=\sup_{\varphi:\Pr_L(\varphi)\le\alpha}\Pr_H(\varphi)\)
denote the ROC curve of the flow experiment.
Here \(\Pr_j(\varphi)=\mathbb E_j\varphi(X)\).

\Needspace{8\baselineskip}

\protect\phantomsection\label{result-lemma-oa-laplace-extreme}{\textbf{Lemma
(Laplace extremality; analytical).}} \emph{(i) If \(f\in\mathcal C_L\)
has median zero, then \(\Pr_f(|Z|>t)\ge e^{-Lt}\) for every \(t\ge0\).
Equality at a positive \(t\) requires \(f=\ell_b\).
(ii) At every pure order profile, the Laplace flow experiment Blackwell-dominates the experiment from any \(f\in\mathcal C_L\).
This dominance carries to prices if each price is a measurable function of its flow and the Laplace price reveals its flow posterior.
(iii) At a common scale, logistic noise equals Laplace noise plus independent noise:
\(Z_{\mathrm{log}}\overset d=Z_{\mathrm{Lap}}+W\), where
\(W\overset d=\sum_{n=2}^{\infty}Y_n\) and the \(Y_n\) are independent Laplace variables with scales \(b/n\).
The flow comparison holds at mixed order profiles as well.
It carries to prices under the same hypotheses.}

\emph{Proof of (i).} Put \(S(t)=\int_t^\infty f(z)\,dz\) for \(t\ge0\).
The log-density bound gives \(f(z)\ge f(t)e^{-L(z-t)}\) for \(z\ge t\).
Hence \(S(t)\ge f(t)/L\).
Thus \(S'(t)=-f(t)\ge-LS(t)\).
With \(S(0)=1/2\), this gives \(S(t)\ge e^{-Lt}/2\).
Apply the same argument to \(f(-z)\) and add the two tail bounds.
\par
If their sum equals \(e^{-Lt}\) at a positive \(t\), equality holds in both tail bounds.
Monotonicity of \(e^{Lu}S(u)\) gives \(S(u)=e^{-Lu}/2\) on \([0,t]\).
Differentiation gives \(f(u)=(L/2)e^{-Lu}\) almost everywhere on \((0,t)\).
Continuity extends this identity to every point of \([0,t]\).
In particular, \(f(t)=Le^{-Lt}/2\), so \(S(t)=f(t)/L\).
\par
The continuous function \(f(z)-f(t)e^{-L(z-t)}\) is nonnegative on \([t,\infty)\).
Its integral is zero, so it vanishes at every point there.
Thus \(f(z)=(L/2)e^{-Lz}\) for every \(z\ge0\).
Reflection fixes the left tail. \(\square\)

\emph{Proof of (ii).} Equal orders give constant information.
If \(q_H<q_L\), reflect the observation and the noise density.
Reflection preserves the score bound and the Laplace law.
It therefore suffices to prove the claim for \(q_L<q_H\).
\par
For a statistical decision rule \(0\le\varphi\le1\), put \(h_\varphi(s)=\int\varphi(x)f(x-s)\,dx\).
Section A.3 gives \(|h_\varphi'|\le Lh_\varphi\).
Applied to \(1-\varphi\), it gives \(|h_\varphi'|\le L(1-h_\varphi)\) as well.
Let \(v'=\Phi(v)=L\min\{v,1-v\}\), with \(v(q_L)=h_\varphi(q_L)=\alpha\).
Since \(\Phi\) is Lipschitz, \((h_\varphi-v)_+'\le L(h_\varphi-v)_+\) almost everywhere.
Integration gives \(h_\varphi\le v\).
\par
For Laplace noise, a half-line rule satisfies \(h'=\Phi(h)\) exactly.
Such rules attain every \(\alpha\in(0,1)\).
\par
Half-lines are optimal by monotone likelihood ratio.
The ROC curves therefore satisfy \(\beta_f(\alpha)\le\beta_{\mathrm{Lap}}(\alpha)\).
Continuity gives the endpoints.
For \(0\le t\le1\), Bayes' rule gives \(\mathbb E(\mu-t)_+=\sup_{0\le\alpha\le1}\{\pi(1-t)\beta(\alpha)-(1-\pi)t\alpha\}\).
The ROC inequality orders every posterior call value.
Both posterior means equal \(\pi\), so the order is convex order.
\par
With two states, convex order is Blackwell dominance.
A price that reveals its flow posterior is Blackwell-equivalent to that flow experiment.
Any measurable function of the other flow is a garbling of that experiment.
The stated price comparison follows by transitivity. \(\square\)

\emph{Proof of (iii).} For the logistic density in (A.21), Euler's product gives the characteristic function
\(\chi_{\mathrm{log}}(t)=\prod_{n=1}^{\infty}(1+b^2t^2/n^2)^{-1}\).
Each factor is the characteristic function of a centered Laplace variable with scale \(b/n\).
Take independent \(Y_n\), \(n\ge1\), with these laws.
Their total variance is \(2b^2\sum_{n\ge1}n^{-2}<\infty\).
The partial sums converge almost surely and in \(L^2\).
Their characteristic functions converge to the product.
\par
Separate the first term.
\(Z_{\mathrm{Lap}}=Y_1\) is independent of \(W=\sum_{n\ge2}Y_n\).
The product identity gives \(Z_{\mathrm{log}}\overset d=Z_{\mathrm{Lap}}+W\).
The kernel that adds \(W\) does not depend on the class or on the investor's order.
It gives Blackwell dominance at every pure or mixed order profile.
The same transitivity argument proves the conditional price comparison. \(\square\)

\emph{Price construction.}
The model satisfies the price hypotheses at fixed orders under the following maintained conditions.
The pair \((j,\theta)\), the incumbent value, the preparation cost, and noise are mutually independent.
Conditional on its class, the challenger value is independent of flow.
The auction objects satisfy \(t_L-t_0\ge0\), \(\Delta_T>0\), and \(g_H\ge g_L\ge0\).
Costs lie in \(\{0,c_H\}\), with probability \(\rho>0\) on zero cost and entry at equality.
Then
\[
e_r(\mu)=\rho+(1-\rho)\mathbf1\{B_r(\mu)\ge c_H\},\qquad
P_r(\mu)=t_0+e_r(\mu)[t_L-t_0+\Delta_T\mu].
\]
Entry is nondecreasing in \(\mu\), and for \(\mu_2>\mu_1\),
\[
P_r(\mu_2)-P_r(\mu_1)\ge\rho\Delta_T(\mu_2-\mu_1)>0.
\]
Thus \(P_r\) has a measurable inverse on its image.
For each noise law, apply \(P_r\) to its flow posterior.
The resulting price is measurable in flow and reveals that posterior.
\par
A flow posterior is a sufficient statistic for its two-class experiment.
Bayes' rule makes the conditional flow law, given its posterior, independent of the class.
These constructed prices therefore satisfy the hypotheses of both comparisons.
This construction does not require the fixed order profile to be an equilibrium.

Part (ii) asserts no mixed-profile dominance for a general member of the class.
Part (iii) gives the mixed-profile comparison for logistic noise.
Blackwell dominance orders the value of information.
It does not order the probability of expensive preparation.
